% GENERATED FILE. Edit canonical lessons, then run npm run generate:books.
% canonical-source-hash: fnv1a64:2fe246990c3caf7d
% canonical-lessons: ZH-C360-gonggong, ZH-C360-biyao, ZH-W360-bi-must, ZH-C360-bubi, ZH-C360-changhe, ZH-C360-chadui

\chapter{No Need}
\label{ch:zh-no-need}
\hlchaptermodality{360}

\begin{chapteropening}
\textbf{By the end of this chapter:} \emph{I can say what is necessary, what there is no need for, and what suits the occasion.}

Complete the last lesson of chapter 360: I can say what is necessary, what there is no need for, and what suits the occasion.
\end{chapteropening}

\section[gōnggòng]{gōnggòng — public; shared by everyone}

\label{lesson:ZH-C360-gonggong}



\begin{quote}
Before the new one: say the Mandarin for to put a limit on, then the Mandarin for morality.
\end{quote}

\subsection*{You'll want to know: gōnggòng}
\emph{gōnggòng} — \textquotedblleft{}public; shared by everyone\textquotedblright{}.

\begin{quote}
\emph{zhèr yǒu gōnggòng cèsuǒ ma} — Is there a public toilet here?
\end{quote}

Its characters come later; for now it is a word for the ear.

\begin{cousinweb}
\emph{gōnggòng} is built from \zh{公} and \emph{gòng}, both inherited from Middle Chinese, \emph{kuwng} and \emph{gjowngH} (Guangyun, in Baxter's transcription), which glosses the first 'shared' and 'upright', and the second 'together' and 'all'.
\end{cousinweb}

\begin{grammarlens}[title={What you've built}]
One more word for giving reasons.
\end{grammarlens}

\subsection*{Your turn}
\begin{itemize}
\raggedright
  \item \emph{Say it:} \emph{gōnggòng}
  \item \emph{Say it:} \emph{gōnggòng}, once more
  \item \emph{Say it:} \emph{zhèr yǒu gōnggòng cèsuǒ ma} — Is there a public toilet here?
  \item \emph{Recall:} say the Mandarin for to put a limit on, then the Mandarin for morality, then say \emph{gōnggòng} again
\end{itemize}

\begin{tcolorbox}[breakable,colback=teal!4,colframe=teal!35!black,title={Before you move on}]
What does gōnggòng mean? (\textquotedblleft{}Public; shared by everyone\textquotedblright{}.) Say it once more.
\end{tcolorbox}
\section[\texorpdfstring{\zh{必要} (bìyào)}{必要 (bìyào)}]{\zh{必要} (bìyào) — necessary}

\label{lesson:ZH-C360-biyao}



\begin{quote}
Before the new one: say the Mandarin for morality, then the Mandarin for shared by everyone.
\end{quote}

\subsection*{You'll want to know: \zh{必要}}
\textbf{\zh{必要}} — \emph{bìyào} — \textquotedblleft{}necessary\textquotedblright{}.

\begin{quote}
\textbf{\zh{我觉得没有必要问这个问题}} \emph{wǒ juéde méiyǒu bìyào wèn zhège wèntí} — I don't think there's any need to ask that.
\end{quote}

The new character, in pieces: \zh{必} — five strokes: a dot on the left, the long curve along the foot that hooks up, as in \zh{心}, a dot on top, a long left-falling stroke down through the curve, and a dot on the right. The next lesson puts it in your hand.

\begin{cousinweb}
\textbf{\zh{必要}} (\emph{bìyào}) is built from \emph{bì} and \zh{要} 'to need', both inherited from Middle Chinese, \emph{pjit} and \emph{'jiewH} (Guangyun, in Baxter's transcription), which glosses the first 'certain' and 'so', and the second 'to bind'.
\end{cousinweb}

\begin{grammarlens}[title={What you've built}]
One more word for giving reasons.
\end{grammarlens}

\subsection*{Your turn}
\begin{itemize}
\raggedright
  \item \emph{Say it:} \emph{bìyào}
  \item \emph{Say it:} \emph{bìyào}, once more
  \item \emph{Say it:} \emph{wǒ juéde méiyǒu bìyào wèn zhège wèntí} — I don't think there's any need to ask that.
  \item \emph{Recall:} say the Mandarin for morality, then the Mandarin for shared by everyone, then say \emph{bìyào} again
\end{itemize}

\begin{tcolorbox}[breakable,colback=teal!4,colframe=teal!35!black,title={Before you move on}]
What does \zh{必要} mean? (\textquotedblleft{}Necessary\textquotedblright{}.) Say it once more.
\end{tcolorbox}
\section[\texorpdfstring{\zh{必} (bì)}{必 (bì)}]{\zh{必} — the strokes of the heart, with a sweep through them}

\label{lesson:ZH-W360-bi-must}



\begin{quote}
The word first, then the character inside it.

\begin{itemize}
\raggedright
  \item \emph{Recall:} say \emph{bìyào} — \textbf{R1}, one lesson back, before it decays
\end{itemize}

\begin{quote}
\zh{心} \zh{要}
\end{quote}

\emph{Read it:} each one aloud, with nothing said first — \emph{xīn}, \emph{yào}
\end{quote}

\subsection*{Script — the strokes of the heart, with a sweep through them}
\hlblockfigure{\detokenize{figures/ZH-W360-bi-must-filmstrip.pdf}}{How \zh{必} is written, stroke by stroke}

\begin{quote}
\zh{必}
\end{quote}

\textbf{One piece, five strokes.} It looks like \zh{心} crossed by a long sweep; the order is its own.

\textbf{A dot on the left}, falling down to the left.

\textbf{The long curve along the foot}, sweeping right and hooking up, as in \zh{心}.

\textbf{A dot on top}, falling to the right.

\textbf{A long left-falling stroke} from the top right down through the curve to the bottom left. Only then:

\textbf{the dot on the right}, last.

Five strokes, and the sweep comes before the last dot.

\subsection*{Your turn}
\begin{itemize}
\raggedright
  \item \emph{Write it:} \zh{必}, all five, with the model in front of you
  \item \emph{Cover:} write it again from memory, and count the pieces, not the strokes
\end{itemize}

\begin{tcolorbox}[breakable,colback=teal!4,colframe=teal!35!black,title={Before you move on}]
How many strokes in \zh{必}? (\textbf{Five}.) Which character you know has the same curve? (\textbf{\zh{心}}.) What comes right before the last dot? (\textbf{The long left-falling stroke}.)

Source: \href{https://raw.githubusercontent.com/chanind/hanzi-writer-data/68d10a4b21150cae5e1ebbd223eed289cf32d90c/data/\%E5\%BF\%85.json}{Hanzi Writer Data \zh{必}, pinned stroke paths and medians}.
\end{tcolorbox}
\section[\texorpdfstring{\zh{不必} (búbì)}{不必 (búbì)}]{\zh{不必} (búbì) — need not; there's no need to}

\label{lesson:ZH-C360-bubi}



\begin{quote}
Before the new one: say the Mandarin for shared by everyone, then the Mandarin for necessary.
\end{quote}

\subsection*{You'll want to know: \zh{不必}}
\textbf{\zh{不必}} — \emph{búbì} — \textquotedblleft{}need not; there's no need to\textquotedblright{}.

\begin{quote}
\textbf{\zh{你不必说对不起你没有错}} \emph{nǐ búbì shuō duìbuqǐ nǐ méiyǒu cuò} — You needn't apologize; you did nothing wrong.
\end{quote}

The tone of bù changes here: before a falling tone it is said bú, and the spelling writes the change.

\begin{cousinweb}
\textbf{\zh{不必}} (\emph{búbì}) is built from \zh{不} 'not' and \zh{必} 'surely', both inherited from Middle Chinese, \emph{pjut} and \emph{pjit} (Guangyun, in Baxter's transcription), which glosses the second 'certain' and 'so'.
\end{cousinweb}

\begin{grammarlens}[title={What you've built}]
One more word for giving reasons.
\end{grammarlens}

\subsection*{Your turn}
\begin{itemize}
\raggedright
  \item \emph{Say it:} \emph{búbì}
  \item \emph{Say it:} \emph{búbì}, once more
  \item \emph{Say it:} \emph{nǐ búbì shuō duìbuqǐ nǐ méiyǒu cuò} — You needn't apologize; you did nothing wrong.
  \item \emph{Recall:} say the Mandarin for shared by everyone, then the Mandarin for necessary, then say \emph{búbì} again
\end{itemize}

\begin{tcolorbox}[breakable,colback=teal!4,colframe=teal!35!black,title={Before you move on}]
What does \zh{不必} mean? (\textquotedblleft{}Need not; there's no need to\textquotedblright{}.) Say it once more.
\end{tcolorbox}
\section[chǎnghé]{chǎnghé — occasion; setting}

\label{lesson:ZH-C360-changhe}



\begin{quote}
Before the new one: say the Mandarin for necessary, then the Mandarin for there's no need to.
\end{quote}

\subsection*{You'll want to know: chǎnghé}
\emph{chǎnghé} — \textquotedblleft{}occasion; setting\textquotedblright{}.

\begin{quote}
\emph{nǐ yīnggāi kàn chǎnghé shuōhuà} — You should suit what you say to the occasion.
\end{quote}

Its characters come later; for now it is a word for the ear.

\begin{cousinweb}
\emph{chǎnghé} is built from \emph{chǎng} and \emph{hé}, both inherited from Middle Chinese, \emph{drjang} and \emph{hop} (Guangyun, in Baxter's transcription), which glosses the first 'a place of offering' and 'a threshing floor', and the second 'to join'. Middle Chinese \emph{drjang} would regularly give cháng; the dipping tone of chǎng is irregular.
\end{cousinweb}

\begin{grammarlens}[title={What you've built}]
One more word for giving reasons.
\end{grammarlens}

\subsection*{Your turn}
\begin{itemize}
\raggedright
  \item \emph{Say it:} \emph{chǎnghé}
  \item \emph{Say it:} \emph{chǎnghé}, once more
  \item \emph{Say it:} \emph{nǐ yīnggāi kàn chǎnghé shuōhuà} — You should suit what you say to the occasion.
  \item \emph{Recall:} say the Mandarin for necessary, then the Mandarin for there's no need to, then say \emph{chǎnghé} again
\end{itemize}

\begin{tcolorbox}[breakable,colback=teal!4,colframe=teal!35!black,title={Before you move on}]
What does chǎnghé mean? (\textquotedblleft{}Occasion; setting\textquotedblright{}.) Say it once more.
\end{tcolorbox}
\section[chāduì]{chāduì — to jump the queue}

\label{lesson:ZH-C360-chadui}



\begin{quote}
Before the new one: say the Mandarin for there's no need to, then the Mandarin for the occasion you are in.
\end{quote}

\subsection*{You'll want to know: chāduì}
\emph{chāduì} — \textquotedblleft{}to jump the queue\textquotedblright{}.

\begin{quote}
\emph{dàjiā dōu zài páiduì nǐ bù yīnggāi chāduì} — Everyone's queuing; you shouldn't push in.
\end{quote}

Its characters come later; for now it is a word for the ear.

\begin{cousinweb}
\emph{chāduì} is built from \emph{chā} and \emph{duì}, both inherited from Middle Chinese, \emph{tsrheap} and \emph{dwojH} (Guangyun, in Baxter's transcription), which glosses the first 'to thrust in' and the second 'a troop'.
\end{cousinweb}

\begin{grammarlens}[title={What you've built}]
One more word for giving reasons.
\end{grammarlens}

\subsection*{Your turn}
\begin{itemize}
\raggedright
  \item \emph{Say it:} \emph{chāduì}
  \item \emph{Say it:} \emph{chāduì}, once more
  \item \emph{Say it:} \emph{dàjiā dōu zài páiduì nǐ bù yīnggāi chāduì} — Everyone's queuing; you shouldn't push in.
  \item \emph{Recall:} say the Mandarin for there's no need to, then the Mandarin for the occasion you are in, then say \emph{chāduì} again
\end{itemize}

\begin{tcolorbox}[breakable,colback=teal!4,colframe=teal!35!black,title={Before you move on}]
What does chāduì mean? (\textquotedblleft{}To jump the queue\textquotedblright{}.) Say it once more.
\end{tcolorbox}
